\chapter{Case Studies}
\label{ch:cases}

The algorithm is tested on case studies of increasing size and difficulty. They differ in the
body (Eros, 1\,499 facets; Kleopatra, 46 times the surface and 4\,092 facets), the longest allowed
mission $T_{\max}$ (from two hours to ten days, i.e.\ from 13 to 1\,441 simulated steps), and the
number of slots $N_{\max}$ (5 or 10, i.e.\ 26 or 51 decision variables). Short missions are the
hard ones: the asteroid turns only a few times, so \emph{where} each spacecraft is \emph{when}
decides what can be seen, and the fitness landscape becomes rugged.

In every case the mission duration is searched in $[T_{\min},T_{\max}]$ with
$T_{\min}=\min(6\,\mathrm{h},\,T_{\max}/4)$, and all other settings are those of
Tables~\ref{tab:params} and~\ref{tab:saparams}. Every \ac{SA} run uses 2000 evaluations, and every
case is repeated with independent random seeds (a seed fixes all random choices of one run:
start designs, moves and acceptances), so that the mean and spread of the results can be
reported. Table~\ref{tab:cases} lists the cases and their input parameters.

\begin{table}[htbp]
\centering
\caption{Case studies and their input parameters (budget 2000 evaluations per run)}
\label{tab:cases}
\footnotesize
\begin{tabularx}{\textwidth}{@{}p{0.5cm}p{2.9cm}p{0.9cm}p{0.9cm}p{0.8cm}p{0.8cm}p{0.6cm}X@{}}
\toprule
Case & Body, setting & $T_{\max}$ & $T_{\min}$ & $N_{\max}$ & dim & Seeds & Question \\
\midrule
A1 & Eros & 10 d & 6 h & 5 & 26 & 5 & long mission: coverage saturates \\
A2 & Eros & 1 d & 6 h & 5 & 26 & 5 & \multirow{5}{=}{How do fleet size, orbits and duration change as the time available shrinks?} \\
A3 & Eros & 1 d & 6 h & 10 & 51 & 5 & \\
A4 & Eros & 12 h & 3 h & 10 & 51 & 5 & \\
A5 & Eros & 6 h & 1.5 h & 10 & 51 & 5 & \\
A6 & Eros & 2 h & 0.5 h & 10 & 51 & 5 & \\
\midrule
B1 & Kleopatra & 3 d & 6 h & 5 & 26 & 5 & \multirow{3}{=}{Does the method carry over unchanged to a 6.5$\times$ larger, strongly concave body?} \\
B2 & Kleopatra & 3 d & 6 h & 10 & 51 & 5 & \\
B3 & Kleopatra & 12 h & 3 h & 10 & 51 & 5 & \\
\midrule
C1 & Eros, $N$ forced to 1--6 & 2 d & 6 h & $N$ & $1{+}5N$ & 5 & \multirow{3}{=}{Is the fleet size \ac{SA} chooses really the best? (V6)} \\
C2 & Eros, $N$ forced to 1--8 & 6 h & 1.5 h & $N$ & $1{+}5N$ & 5 & \\
C3 & Kleopatra, $N$ forced to 1--5 & 3 d & 6 h & $N$ & $1{+}5N$ & 3 & \\
\midrule
D & Eros, no Sun / shadows & 10 d & 6 h & 5 & 26 & 5 & comparison with the literature (V5) \\
E & Eros, SA parameters varied & 2 d & 6 h & 5 & 26 & 5 & effect of $\alpha$, $p_0$, $\sigma_0$, warm start \\
\bottomrule
\end{tabularx}
\end{table}

Cases A and B are also solved by random search with the same budget (validation V4). Cases C1--C3
fix every switch to 1 (so $N=N_{\max}$) and optimise only the orbits and the duration; each is
compared with a free-$N$ \ac{SA} run of the same mission with $N_{\max}=10$. In case~E one
parameter at a time is changed from the defaults of Table~\ref{tab:saparams} (values in
Section~\ref{sec:params}).

Table~\ref{tab:cost} shows how the cost of one run grows with the case. The time per evaluation
depends mostly on how many new spacecraft tracks have to be simulated and on their length
($T_{\max}/\Delta t$ steps) times the number of facets.

\begin{table}[htbp]
\centering
\caption{Size of the case studies and the cost of one \ac{SA} run (2000 evaluations, one CPU core)}
\label{tab:cost}
\scriptsize
\begin{tabular}{@{}lllccccc@{}}
\toprule
Setting & Body & $T_{\max}$ & $N_{\max}$ & Variables & Steps $N_t$ & Minutes per run & Evaluations / s \\
\midrule
Eros 10 d & Eros & 10 d & 5 & 26 & 1441 & 2.9 & 11.3 \\
Eros 1 d & Eros & 1 d & 5 & 26 & 145 & 0.3 & 116.0 \\
Eros 1 d & Eros & 1 d & 10 & 51 & 145 & 0.3 & 119.1 \\
Eros 12 h & Eros & 12 h & 10 & 51 & 73 & 0.2 & 187.6 \\
Eros 6 h & Eros & 6 h & 10 & 51 & 37 & 0.1 & 332.0 \\
Eros 2 h & Eros & 2 h & 10 & 51 & 13 & 0.1 & 567.7 \\
Kleopatra 3 d & Kleopatra & 3 d & 5 & 26 & 433 & 3.6 & 9.3 \\
Kleopatra 3 d & Kleopatra & 3 d & 10 & 51 & 433 & 3.2 & 10.3 \\
Kleopatra 12 h & Kleopatra & 12 h & 10 & 51 & 73 & 0.6 & 54.2 \\
\bottomrule
\end{tabular}
\end{table}

